\section{Simulation-Based BCA Representation}

The purpose of the simulation exercise is not to evaluate conventional BCA as an
empirical procedure. Instead, the purpose is constructive. Unlike the original
CKM equivalence results for first-moment frictions, a clean analytical
equivalence between second-moment shocks and canonical BCA wedges is hard to
derive. We therefore define an operational object: the BCA representation of
uncertainty shocks.

\begin{definition}
The BCA representation of an uncertainty shock (either TFP or financial) is the
set of canonical prototype wedges that, when estimated in a CKM-style accounting
framework, reproduce the allocations generated by a nonlinear economy with
stochastic volatility.
\end{definition}

\subsection{True Data-Generating Model}

The simulation uses a representative-agent RBC economy augmented with a unified
collateral constraint. Following the methodological advice in the comprehensive
review by Fern\'andez-Villaverde and Guerr\'on-Quintana (2020), we capture the
essence of financial frictions using a reduced-form collateral constraint within
a representative-agent framework. As they demonstrate, uncertainty in financial
frictions is crucial: without it, uncertainty shocks may counterfactually stimulate
the economy via the Oi-Hartman-Abel effect. The stochastic financial constraints
ensure that uncertainty correctly manifests as massive endogenous risk premiums,
driving the necessary contractionary comovements in the BCA wedges. Specifically,
following Liu, Wang, and Zha (2013) and Chatterjee, Gunawan, and Kohn (2024),
we assume that firms face a credit limit that simultaneously restricts both their
intertemporal borrowing for capital investment and their intratemporal borrowing
for working capital (i.e., financing the wage bill).

The firm chooses labor, investment, and debt to maximize expected discounted
profits, evaluated using the household's stochastic discount factor \(\Lambda_{0,t}\):
\begin{equation}
    \max_{\{L_t, I_t, K_{t+1}, B_t\}} \E_0 \sum_{t=0}^{\infty} \Lambda_{0,t}
    \left[ Y_t - W_t L_t - I_t + B_t - \frac{R_{t-1}}{\Pi_t}B_{t-1} \right],
    \label{eq:simulation_firm_problem}
\end{equation}
subject to the production function \(Y_t = A_t K_t^\alpha L_t^{1-\alpha}\), the
capital accumulation equation \(K_{t+1}=(1-\delta)K_t+I_t\), and the unified
collateral constraint:
\begin{equation}
    B_t + \phi_t W_t L_t \leq \theta_t \E_t \left[ \frac{\Pi_{t+1}}{R_t} q_{t+1} K_{t+1} \right].
    \label{eq:simulation_collateral}
\end{equation}
Here, \(B_t\) is intertemporal debt, \(\phi_t \in [0,1]\) is the stochastic fraction
of the wage bill that must be financed in advance via working capital, \(q_{t+1}\)
is the shadow price of capital, \(\Pi_{t+1}/R_t\) is the inverse of the real interest
rate, and \(\theta_t\) is the stochastic loan-to-value (LTV) ratio.

By specifying both \(\theta_t\) and \(\phi_t\) as distinct stochastic processes, we
elegantly nest the two leading financial friction mechanisms in the literature.
The collateral shock \(\theta_t\) (as in Liu, Wang, and Zha 2013) primarily
governs the overall borrowing capacity. Concurrently, the working capital shock
\(\phi_t\) governs the liquidity drain imposed by payrolls. Notably, the pure
working capital constraint in Chatterjee, Gunawan, and Kohn (2024)---where labor
demand is constrained by a stochastic fraction \(\zeta_t\) of capital value---emerges
as a special case when intertemporal debt is absent (\(B_t=0\)), mapping their
credit supply shock \(\zeta_t\) precisely to the ratio \(\theta_t / \phi_t\).

As highlighted by Lin, Tsai, and Wang (2025), formulating the collateral constraint
in terms of the objective real interest rate \(\Pi_{t+1}/R_t\) fundamentally alters
asset pricing. The firm's optimal choice of capital is determined by a hybrid
discount factor that forces the firm to bear substantial covariance risk:
\begin{equation}
    q_t = \E_t \left[ \Lambda_{t+1,t} \alpha \frac{Y_{t+1}}{K_{t+1}} \right]
    + \E_t \left[ \left( (1-\delta)\Lambda_{t+1,t} + \mu_t \theta_t \frac{\Pi_{t+1}}{R_t} \right) q_{t+1} \right],
    \label{eq:simulation_capital_foc}
\end{equation}
where \(\mu_t = 1 - \E_t[\Lambda_{t+1,t} R_t / \Pi_{t+1}]\) is the Lagrange multiplier on
the collateral constraint, representing the external finance premium. Simultaneously,
the working capital requirement distorts labor demand:
\begin{equation}
    (1-\alpha)\frac{Y_t}{L_t} = W_t (1 + \phi_t \mu_t).
    \label{eq:simulation_labor_foc}
\end{equation}
This two-shock structure perfectly separates the transmission into the BCA space.
When uncertainty in the collateral limit (\(\theta_t\)) rises, the endogenous risk
premium \(\mu_t\) explodes, violently distorting the Euler equation \eqref{eq:simulation_capital_foc}
and driving the BCA investment wedge. Conversely, when uncertainty in the working
capital requirement (\(\phi_t\)) rises, it directly and aggressively maps into the
BCA labor wedge \((1-\tau_{\ell,t}) = 1/(1+\phi_t\mu_t)\).

The households maximize lifetime utility:
\begin{equation}
    \max_{\{C_t, L_t\}} \E_0 \sum_{t=0}^{\infty} \beta^t
    \left[ \frac{C_t^{1-\gamma}-1}{1-\gamma} - \chi\frac{L_t^{1+\nu}}{1+\nu} \right],
    \label{eq:simulation_ra_problem}
\end{equation}
subject to their budget constraint \(C_t + \frac{R_{t-1}}{\Pi_t}B_{t-1} = W_t L_t + \Pi_t + B_t\).
Because firms are owned by households, the aggregate resource constraint remains frictionless:
\begin{equation}
    Y_t = C_t + I_t.
    \label{eq:simulation_resource}
\end{equation}

The economy is driven by three independent stochastic volatility processes. The
first governs TFP (Macro) and its uncertainty:
\begin{align}
    \log A_t &= \rho_A \log A_{t-1} + \exp(V_{A,t})\sigma_A\varepsilon_{A,t},
    \label{eq:simulation_tfp} \\
    V_{A,t} &= \rho_{VA} V_{A,t-1} + \sigma_{VA}\varepsilon_{VA,t}.
    \label{eq:simulation_tfp_vol}
\end{align}
The second governs the LTV ratio (Collateral Financial Friction):
\begin{align}
    \log \theta_t &= (1-\rho_\theta)\log \bar{\theta} + \rho_\theta \log \theta_{t-1} + \exp(V_{\theta,t})\sigma_\theta\varepsilon_{\theta,t},
    \label{eq:simulation_theta} \\
    V_{\theta,t} &= \rho_{V\theta} V_{\theta,t-1} + \sigma_{V\theta}\varepsilon_{V\theta,t}.
    \label{eq:simulation_theta_vol}
\end{align}
The third governs the working capital requirement (Liquidity Financial Friction):
\begin{align}
    \log \phi_t &= (1-\rho_\phi)\log \bar{\phi} + \rho_\phi \log \phi_{t-1} + \exp(V_{\phi,t})\sigma_\phi\varepsilon_{\phi,t},
    \label{eq:simulation_phi} \\
    V_{\phi,t} &= \rho_{V\phi} V_{\phi,t-1} + \sigma_{V\phi}\varepsilon_{V\phi,t}.
    \label{eq:simulation_phi_vol}
\end{align}

In the deterministic steady state (DSS), we calibrate \(\bar{\theta}\) and \(\bar{\phi}\)
such that the constraint binds, mapping the model to a standard RBC economy with
a steady-state financial distortion. The financial uncertainty shocks (\(\varepsilon_{V\theta,t}\)
and \(\varepsilon_{V\phi,t}\)) act as pure second-moment shocks that trigger precautionary
labor reductions and investment contractions through the endogenous risk premium,
without directly altering first-moment productivity.

\subsection{Calibration Strategy}

To operationalize the simulation and generate the BCA representation prior to
finalizing the empirical OI-SVMVAR estimates, we adopt a preliminary calibration
strategy firmly grounded in the leading DSGE literature on financial frictions.
This ensures the baseline quantitative results are disciplined by established
macroeconomic moments.

\paragraph{Macroeconomic block and TFP uncertainty.}
The structural parameters for the representative-agent RBC core follow standard
values: \(\beta = 0.99\) (4\% annual real interest rate), \(\alpha = 0.33\) (capital
share), \(\delta = 0.025\) (10\% annual depreciation), and \(\gamma = 1\) (log
utility in consumption). The Frisch elasticity of labor supply is set to 1
(\(\nu = 1\)). The baseline parameter values for the TFP process and its stochastic
volatility match the estimated values in the dissertation manuscript:
\[
    \rho_A=0.9780,\qquad
    \sigma_A=0.0089,\qquad
    \rho_{VA}=0.9110,\qquad
    \sigma_{VA}=0.1643.
\]

\paragraph{Collateral financial friction (\(\theta_t\)).}
We calibrate the collateral limit process following the Bayesian estimates of
Liu, Wang, and Zha (2013). The steady-state loan-to-value ratio is set to
\(\bar{\theta} = 0.75\), matching the average ratio of nonfarm nonfinancial
business debt to tangible assets in the U.S. The persistence of the first-moment
shock is set to \(\rho_\theta = 0.98\), reflecting the highly persistent nature
of credit conditions. The baseline standard deviation of the collateral shock
is \(\sigma_\theta = 0.0112\).

\paragraph{Liquidity financial friction (\(\phi_t\)).}
The working capital requirement process is calibrated using the estimates from
Chatterjee, Gunawan, and Kohn (2024), who explicitly model the stochastic
volatility of credit growth. We set the steady-state fraction of the wage bill
requiring advance financing to \(\bar{\phi} = 1.0\), standard in the working
capital literature (e.g., Christiano, Eichenbaum, and Evans, 2005). The persistence
of the first-moment liquidity shock is \(\rho_\phi = 0.836\), with a baseline
standard deviation of \(\sigma_\phi = 0.01\).

\paragraph{Financial uncertainty processes (\(V_{\theta,t}\) and \(V_{\phi,t}\)).}
Because the ultimate goal is to map the theoretical financial friction shocks to
the empirical "Financial Factor" extracted from the OI-SVMVAR, we assume for this
preliminary simulation that both financial frictions share a common uncertainty
process. Formally, define the common financial volatility state and innovation as
\[
    V_{F,t}\equiv V_{\theta,t}=V_{\phi,t},
    \qquad
    \varepsilon_{VF,t}\equiv
    \varepsilon_{V\theta,t}=\varepsilon_{V\phi,t}.
\]
We draw the stochastic volatility parameters directly from the credit
uncertainty estimates of Chatterjee et al. (2024):
\[
    \rho_{V\theta} = \rho_{V\phi} = 0.867,\qquad
    \sigma_{V\theta} = \sigma_{V\phi} = 0.359.
\]
Notably, the volatility of financial shocks (\(0.359\)) is substantially larger
than the volatility of TFP shocks (\(0.1643\)). This large second-moment variation
is crucial: it triggers the massive endogenous risk premium (\(\mu_t\)) required
to violently distort the BCA labor and investment wedges during a financial
uncertainty episode.

\subsection{Deterministic vs Stochastic Steady State}

The true model is solved by third-order perturbation with pruning. A first-order
solution is not useful because certainty equivalence eliminates the effect of
\(\varepsilon_{VA,t}\) and \(\varepsilon_{VF,t}\) on allocations. Two distinct
steady-state concepts matter for the experimental design.

The \emph{deterministic steady state} (DSS) is the fixed point of the model
under the counterfactual of agents facing \(\sigma_A=\sigma_f=\sigma_{VA}=\sigma_{VF}=0\).
It is computed analytically via the \texttt{steady\_state\_model} block and is
the point around which Dynare expands the policy functions.

The \emph{stochastic steady state} (SSS) is the fixed point of the
\emph{stochastic-environment} policy functions when all realized innovations
are set to zero. Under third-order perturbation, the policy functions carry
precautionary terms proportional to the shock variances, so with
realized shocks at zero the system still drifts away from DSS and settles at a
distinct point SSS. Numerically, SSS is obtained by initializing the solved
policy functions at DSS and iterating forward under a zero innovation sequence
until convergence.

Under a first-order approximation SSS collapses onto DSS by certainty
equivalence; the fact that SSS \(\neq\) DSS in our setting is precisely the
object that makes uncertainty shocks nontrivial.

\subsection{Two-Experiment Design}

Both experiments share a common warm-up to SSS. They differ in how the
identification of the uncertainty shocks is constructed. The updated design
estimates the orthogonal responses to both Macro (TFP) and Financial uncertainty.

\paragraph{Warm-up (shared across experiments).}
Start from DSS. Iterate the third-order pruned policy functions forward for
\(T^{\text{wu}}=50{,}000\) periods with all innovations fixed at zero. The
resulting state vector is the numerical SSS. Convergence is verified by
comparing the state at \(T^{\text{wu}}\) against the state at
\(T^{\text{wu}}-1{,}000\).

\paragraph{Experiment~1: ergodic sample with local projections.}
Starting from SSS, simulate a single long path of length \(T=20{,}000\) with
all four innovations \(\overset{\text{iid}}{\sim}\mathcal{N}(0,1)\).
Apply the CKM estimator described in \S\ref{subsec:ckm_estimator} to recover
\(\{\widehat{s}_t\}_{t=1}^T\) where
\(s_t=(\log A_t,\,G_t/Y_t,\,\tau_{\ell,t},\,\tau_{x,t})'\). Then estimate local
projections:
\begin{equation}
    \widehat{s}_{j,t+h}
    =
    \alpha_{j,h}
    +
    \beta_{j,h}^{\text{macro}}\varepsilon_{VA,t}
    +
    \beta_{j,h}^{\text{fin}}\varepsilon_{VF,t}
    +
    \Gamma_{j,h}'Z_t
    +
    u_{j,t+h},
    \qquad
    j\in\{A,G,\ell,x\},
    \label{eq:lp_sample}
\end{equation}
where \(Z_t\) includes first-moment shocks (\(\varepsilon_{A,t}, \varepsilon_{f,t}\)),
lagged \(V_{A,t}, V_{F,t}\), and lagged recovered wedges. The estimated sequences
\(\{\beta_{j,h}^{\text{macro}}\}\) and \(\{\beta_{j,h}^{\text{fin}}\}\) represent
the time-average response of wedge \(j\) to each uncertainty innovation.

\paragraph{Experiment~2: generalized impulse responses via antithetic Monte Carlo panel.}
From the same SSS, draw \(N\) innovation sequences for all four shocks.
To isolate the financial uncertainty shock, for example, construct two matched paths:
\begin{align}
    \text{baseline:}\quad
    & \varepsilon_{VF,0}^{(i)}=0,
      \quad
      \text{all other shocks drawn as usual},
    \\
    \text{shocked:}\quad
    & \varepsilon_{VF,0}^{(i)}=1,
      \quad
      \text{all other shocks drawn exactly as in baseline}.
\end{align}
Post-shock innovations are identical within each pair, so the only difference
is the period-zero financial uncertainty shock. Apply the CKM estimator to both
paths, subtract, and average across \(N\):
\begin{equation}
    \widehat{\mathrm{GIRF}}^{\text{fin}}_{j,h}
    =
    \frac{1}{N}\sum_{i=1}^{N}
    \Big[
        \widehat{s}^{(i),\text{shock}}_{j,h}
        -
        \widehat{s}^{(i),\text{base}}_{j,h}
    \Big],
    \qquad
    j\in\{A,G,\ell,x\}.
    \label{eq:girf}
\end{equation}
Repeat this procedure symmetrically to construct \(\widehat{\mathrm{GIRF}}^{\text{macro}}_{j,h}\)
by shocking \(\varepsilon_{VA,0}\). This captures both the mean response to the
volatility state and the second-moment amplification of subsequent first-moment
innovations. Suggested scale: \(N=1{,}000\) antithetic pairs and horizon \(H=60\).

\paragraph{Cross-validation.}
Under ergodicity and smoothness of the policy functions, equations
\eqref{eq:lp_sample} and \eqref{eq:girf} describe asymptotically equivalent
objects. Disagreement between \(\beta\) and \(\widehat{\mathrm{GIRF}}\)
measures the magnitude of higher-order nonlinearity in the mapping from uncertainty
to BCA wedges.

\subsection{CKM-Style Wedge Estimation}
\label{subsec:ckm_estimator}

The simulated data \(\{Y_t,C_t,X_t,L_t,K_t\}\) are treated as if they were
ordinary macroeconomic data. The accounting model is the standard prototype
with four canonical wedges:
\[
    A_t,\qquad
    G_t,\qquad
    \tau_{\ell,t},\qquad
    \tau_{x,t}.
\]
The wedges should be estimated in a CKM-style system rather than simply read
off equation by equation. In particular, the investment wedge depends on
expected future marginal utilities and returns, so its recovery should be
disciplined by the estimated wedge law of motion. A minimal implementation
can use a VAR(1) law of motion for transformed wedges:
\begin{equation}
    s_{t+1}
    =
    P_0 + P s_t + \varepsilon_{t+1},
    \qquad
    s_t
    =
    \begin{bmatrix}
        \log A_t \\
        G_t/Y_t \\
        \tau_{\ell,t} \\
        \tau_{x,t}
    \end{bmatrix}.
    \label{eq:ckm_wedge_law}
\end{equation}
The exact state vector and transformations should follow the CKM
implementation as closely as possible in the coding stage. Because the true
DGP has frictionless resource constraint (\(G_t\equiv 0\)), the recovered
\(\widehat{G}_t/Y_t\) series serves as a tripwire: material deviation from
zero signals a resource-identity violation in the simulated data.

\subsection{Implementation Roadmap}

The executable scaffold is placed in
\texttt{code/bca\_uncertainty\_simulation/}. The theory note remains in this
folder while simulation code lives under \texttt{code/} so that it can later
be extended into a reproducible estimation pipeline.

The coding stage proceeds in five steps:
\begin{enumerate}[label=\textbf{Step \arabic*.}]
    \item \textbf{Solve the model.} Run \texttt{stoch\_simul(order=3,
    k\_order\_solver, pruning, periods=0, irf=0)} so that only the
    decision-rule object \texttt{oo\_.dr} is produced. No simulation is
    performed inside Dynare itself.
    \item \textbf{Shared warm-up to SSS.} From DSS, iterate the pruned policy
    functions forward for \(T^{\text{wu}}=50{,}000\) periods with all
    innovations set to zero. Record the terminal state as SSS and verify
    convergence against \(T^{\text{wu}}-1{,}000\).
    \item \textbf{Experiment~1.} From SSS, simulate one path of length
    \(T=20{,}000\) with iid standard-normal innovations. Apply the CKM
    estimator and estimate local projections \eqref{eq:lp_sample} on
    \(\varepsilon_{VA,t}\) and \(\varepsilon_{VF,t}\).
    \item \textbf{Experiment~2.} From SSS, generate \(N=1{,}000\) antithetic
    innovation pairs spanning \(H=60\) periods for both Macro and Financial
    uncertainty shocks. Apply the CKM estimator to each path and compute
    \(\widehat{\mathrm{GIRF}}\) in \eqref{eq:girf}.
    \item \textbf{Cross-validate.} Compare \(\beta\) and
    \(\widehat{\mathrm{GIRF}}\). Report both panels and interpret
    disagreement as evidence of higher-order nonlinearity in the
    uncertainty-to-wedge mapping.
\end{enumerate}

Housing is excluded from the first version. Adding housing would introduce a
second durable asset Euler equation and a housing-investment wedge-like
object, obscuring the baseline mapping from the uncertainty shocks to the classical
four-wedge BCA representation.
